% Generated draft: requires mathematical and source review.
\documentclass[11pt]{article}
\usepackage[margin=25mm]{geometry}
\usepackage{fontspec}
\setmainfont{Latin Modern Roman}
\usepackage{amsmath,amssymb,mathtools,unicode-math}
\setmathfont{Latin Modern Math}
\usepackage{enumitem,needspace,xurl,hyperref}
\setlength{\parindent}{0pt}
\setlength{\parskip}{5pt}
\setlength{\emergencystretch}{3em}
\newcommand{\N}{\mathbb{N}}
\newcommand{\Z}{\mathbb{Z}}
\newcommand{\Q}{\mathbb{Q}}
\newcommand{\R}{\mathbb{R}}
\newcommand{\C}{\mathbb{C}}
\newcommand{\F}{\mathbb{F}}
\newcommand{\A}{\mathbb{A}}
\newcommand{\PP}{\mathbb{P}}
\newcommand{\E}{\mathbb{E}}
\newcommand{\eps}{\varepsilon}
\newcommand{\dd}{\,\mathrm{d}}
\begin{document}
\section*{UnsolvedMath ID: 30006390}
\textbf{Source problem code:} catalogue entry
\subsection{Definitions and mathematical statement}
Let \(q\) range over prime powers tending to infinity, as stipulated in the original conjecture. For each such \(q\), let \(\Pi_q=(P_q,\mathcal L_q)\) be a projective plane of order \(q\): any two distinct points lie on exactly one common line, any two distinct lines meet in exactly one point, and each line contains exactly \(q+1\) points. Thus
\[
|P_q|=|\mathcal L_q|=q^2+q+1.
\]
Choose \(R_q\subseteq P_q\) by including each point independently with probability \(1/2\). Define
\[
\mathcal H_q(R_q)=\{L\cap R_q:L\in\mathcal L_q,\ L\cap R_q\ne\varnothing\}
\]
and its transversal number
\[
\tau_q(R_q)=\min\bigl\{|B|:B\subseteq R_q,\ B\cap H\ne\varnothing\text{ for every }H\in\mathcal H_q(R_q)\bigr\}.
\]
In particular, \(\tau_q(\varnothing)=0\). All probabilities below concern only the independent choices defining \(R_q\).

The supplied Conjecture 5 asks whether there exists a deterministic function \(f\), independent of the random outcome, such that
\[
f(q)\longrightarrow\infty,
\qquad
\Pr\!\left(\tau_q(R_q)>qf(q)\right)\longrightarrow1.
\]
Equivalently, with probability tending to one, every \(B\subseteq R_q\) meeting every nonempty line section satisfies \(|B|/q>f(q)\). The proposed stronger assertion asks whether one can choose \(f(q)=\Omega(\log q)\); equivalently, whether some constant \(c>0\) satisfies
\[
\Pr\!\left(\tau_q(R_q)>c q\log q\right)\longrightarrow1,
\]
where \(\log\) denotes the natural logarithm.

The original excerpt requires intersection with all line sections, whereas the cleaned statement explicitly excludes empty sections. These requirements agree on the event that every line meets \(R_q\). Its complement has probability at most
\[
(q^2+q+1)2^{-(q+1)}\longrightarrow0.
\]
Thus this distinction does not change the intended high-probability assertion, although the original minimum is undefined when an empty section occurs unless a convention is added.

For precision, the displayed limits can be read for a prescribed family \((\Pi_q)_q\). The source says only ``a projective plane'' and does not explicitly specify uniformity over choices of planes or whether the same function \(f\), or constant \(c\), must work for all such families. That quantifier qualification remains unresolved. The prime-power restriction above follows the original excerpt; extending the assertion to every order admitting a projective plane is not explicitly justified by that excerpt.
\subsection{Short English statement}
Independently retain each point of a projective plane of prime-power order \(q\) with probability \(1/2\). The conjecture asks whether, with probability tending to one, every subset of retained points meeting every nonempty retained line section has size greater than \(qf(q)\) for some deterministic \(f(q)\to\infty\). A stronger proposed bound is \(c q\log q\) for some \(c>0\). Uniformity over different projective planes remains unspecified in the supplied source.
\subsection{Sources}
\begin{itemize}
\item Supplied local UnsolvedMath registry record, permanent ID \(30006390\), title: \emph{Transversals in Random Subsets of Projective Planes}; collection: Oberwolfach Reports Open Problems; exact source problem code: \texttt{catalogue entry}. The original excerpt identifies the target as Conjecture 5. No individual proposer is supplied; the proposed year is recorded as 2025.
\item Supplied source citation: \emph{Combinatorics, Probability and Computing (2026)}, Oberwolfach Reports, pp.\ 2243--2322, DOI \texttt{10.4171/owr/2025/42}. \url{https://doi.org/10.4171/owr/2025/42}. The citation's 2026 designation and the report's 2025 identifiers are retained as supplied.
\item Supplied literature source: \emph{Oberwolfach Report 42 2025}. \url{https://oa.tib.eu/renate/bitstreams/72e281a6-ba33-4ac7-baf7-6fadd1126345/download}.
\item Supplied literature source, labeled \emph{Alon Blocking Partial Designs and Block-Compatible Sequences}. \url{https://web.math.princeton.edu/~nalon/PDFS/remark191.pdf}.
\item Literature assessment dated \texttt{2026-08-22}: the supplied curation labels the conjecture open and reports that related container results have substantially different parameters and do not establish the requested growing lower bound for the transversal subset. This is supplied dated curation derived from the Oberwolfach report, not an independently verified assessment. The record explicitly requests independent verification of time-sensitive status and cited claims.
\end{itemize}
\end{document}
